% !TEX program = xelatex
% ============================================================
% Подготовка к защите РГР (вариант 16).
% Номера формул берутся из отчёта: сначала собрать ../report/main.tex,
% затем этот файл (latexmk defense.tex).
% ============================================================
\documentclass[a4paper,12pt]{article}

\usepackage{fontspec}
\setmainfont{Times New Roman}
\usepackage{polyglossia}
\setmainlanguage{russian}
\setotherlanguage{english}
\setmonofont{Courier New}
\newfontfamily\cyrillicfonttt{Courier New}

\usepackage[left=22mm,right=18mm,top=18mm,bottom=20mm]{geometry}
\usepackage{setspace}
\setstretch{1.12}
\setlength{\parindent}{0pt}
\setlength{\parskip}{4pt plus 1pt}
\setlength{\emergencystretch}{2em}

\usepackage{amsmath,amssymb}
\renewcommand{\le}{\leqslant}
\renewcommand{\ge}{\geqslant}
\usepackage{array,longtable,booktabs}
\usepackage{enumitem}
\setlist{itemsep=1pt,topsep=2pt,leftmargin=1.6em}
\usepackage{titlesec}
\titleformat{\section}{\normalfont\large\bfseries}{\thesection}{0.6em}{}
\titleformat{\subsection}{\normalfont\normalsize\bfseries}{\thesubsection}{0.6em}{}
\titlespacing*{\section}{0pt}{14pt}{6pt}
\titlespacing*{\subsection}{0pt}{10pt}{4pt}

% Номера формул отчёта: \rf{eq:laplace} -> (9)
\usepackage{xr-hyper}
\usepackage[unicode,hidelinks]{hyperref}
\externaldocument[R-]{../report/main}
\newcommand{\rf}[1]{(\ref*{R-#1})}

\newcommand{\pd}[2]{\dfrac{\partial #1}{\partial #2}}
\newcommand{\pdd}[2]{\dfrac{\partial^2 #1}{\partial #2^2}}
\DeclareMathOperator{\Div}{div}
\DeclareMathOperator{\Grad}{grad}

% Вопрос — ответ ([4] у \nopagebreak обязателен: иначе ответ, начинающийся
% с квадратной скобки, теряет её содержимое)
\newcounter{qa}
\newcommand{\q}[1]{\par\medskip\refstepcounter{qa}%
  \textbf{\theqa.~#1}\par\nopagebreak[4]}
% Разбор формулы
\newcommand{\fhead}[2]{\par\medskip\textbf{Формула #1 — #2.}\par\nopagebreak[4]}

\begin{document}

\begin{center}
  {\Large\bfseries Подготовка к защите РГР, вариант 16\par}
  \vspace{4pt}
  {\large «Управление распределёнными системами».\\
   Этап 1: постановка задачи и математическая модель\par}
\end{center}

\textbf{Как пользоваться.} Номера формул в скобках — это номера из отчёта.
Порядок подготовки: 1) выучить рассказ из раздела~\ref{sec:story} так, чтобы
произносить его своими словами; 2) по разделу~\ref{sec:formulas} пройти
каждую формулу и оба рисунка отчёта: что означает, откуда взялось;
3) научиться воспроизводить на бумаге выводы из раздела~\ref{sec:derivations};
4) прогнать вопросы из раздела~\ref{sec:qa}, закрывая ответы. В конце —
список ловушек и лист самопроверки.

\tableofcontents

% ============================================================
\section{Рассказ на три минуты}
\label{sec:story}

\textbf{Что за задача.} Нужно найти установившуюся температуру в сплошном
цилиндре радиуса~$a$ и высоты~$h$. Торцы теплоизолированы. На боковой
поверхности температура задана: ноль ниже уровня $z=c$ и $T_0$ выше него.

\textbf{Откуда уравнение.} Тепло в твёрдом теле переносится
теплопроводностью. По закону Фурье плотность потока $\vec q=-k\,\Grad T$.
Записываю баланс тепла для произвольного объёма: скорость изменения запаса
тепла равна притоку через поверхность плюс тепло источников. По формуле
Остроградского превращаю поверхностный интеграл в объёмный. Объём
произволен, значит, подынтегральная функция равна нулю — это уравнение
теплопроводности $C\rho\,T_t=\Div(k\,\Grad T)+F$. Материал однородный —
$k$ постоянен. Режим стационарный — $T_t=0$. Источников нет — $F=0$.
Остаётся уравнение Лапласа $\Delta T=0$.

\textbf{Почему цилиндрические координаты.} В координатах $r,\varphi,z$
граница цилиндра состоит из координатных поверхностей $r=a$, $z=0$, $z=h$.
Оператор Лапласа в этих координатах содержит три слагаемых. Цилиндр — тело
вращения, условия на поверхности от угла не зависят, поэтому и температура
от угла не зависит. Слагаемое с $\varphi$ пропадает, остаётся
$T_{rr}+\frac1r T_r+T_{zz}=0$ в прямоугольнике $0<r<a$, $0<z<h$.

\textbf{Граничные условия.} На боковой поверхности задана сама
температура — условие I~рода (Дирихле): $T(a,z)=\mu(z)$. На торцах
теплоизоляция: поток через торец равен нулю, а поток равен
$-k\,\partial T/\partial n$, поэтому $\partial T/\partial z=0$ при $z=0$
и $z=h$ — однородные условия II~рода (Неймана). Ось $r=0$ — не граница
тела, но коэффициент $1/r$ там неограничен, и у уравнения есть решения вида
$\ln r$, поэтому требую ограниченности температуры на оси.

\textbf{Что это за задача.} Уравнение линейное, однородное, второго
порядка, эллиптического типа: $a_{11}=a_{22}=1$, $a_{12}=0$,
$D=a_{12}^2-a_{11}a_{22}=-1<0$. Задача краевая: времени и начальных условий
нет. Решение единственно — это следует из формулы Остроградского для
разности двух решений; из единственности строго получается и независимость
от угла. По принципу максимума $0\le T\le T_0$. В безразмерных переменных
задача зависит только от двух чисел: $a/h$ и $c/h$, а $T_0$ — просто
масштаб.

% ============================================================
\section{Обозначения}
\label{sec:notation}

\begin{longtable}{@{}>{\raggedright\arraybackslash}p{2.6cm}>{\raggedright\arraybackslash}p{9.6cm}>{\raggedright\arraybackslash}p{3.6cm}@{}}
\toprule
Символ & Смысл & Единицы \\
\midrule
\endhead
$T$, $T(r,z)$ & температура в точке цилиндра — искомая функция & К \\
$T_0$ & заданная температура верхней части боковой поверхности (разность температур двух частей); как и в отчёте, для определённости $T_0>0$ & К \\
$T_{\text{н}}$, $T_{\text{в}}$ & температуры нижней и верхней частей боковой поверхности в общем случае (раздел~1) & К \\
$a$ & радиус цилиндра & м \\
$h$ & высота цилиндра & м \\
$c$ & высота линии скачка температуры на боковой поверхности, $0<c<h$ & м \\
$r$ & расстояние от точки до оси цилиндра & м \\
$\varphi$ & полярный угол вокруг оси, $0\le\varphi<2\pi$ & рад \\
$z$ & координата вдоль оси, от нижнего торца & м \\
$x$, $y$ & декартовы координаты в плоскости торца & м \\
$O$, $Oz$ & начало координат (центр нижнего торца) и ось цилиндра & — \\
$t$ & время (нужно только при выводе уравнения) & с \\
$M$ & точка тела & — \\
$\vec q$ & вектор плотности теплового потока & Вт/м$^2$ \\
$q_n$ & проекция $\vec q$ на внешнюю нормаль: тепло через единицу площади поверхности за единицу времени & Вт/м$^2$ \\
$k$ & коэффициент теплопроводности & Вт/(м$\cdot$К) \\
$C$ & удельная теплоёмкость (в лекциях — $c$) & Дж/(кг$\cdot$К) \\
$\rho$ & плотность & кг/м$^3$ \\
$F$ & мощность источников тепла в единице объёма & Вт/м$^3$ \\
$\chi=k/(C\rho)$ & коэффициент температуропроводности (в лекциях — $a^2$) & м$^2$/с \\
$V$, $S$ & произвольный объём внутри тела и его поверхность & м$^3$, м$^2$ \\
$S_\Omega$ & поверхность цилиндра: боковая поверхность и два торца & м$^2$ \\
$\vec n$ & единичная внешняя нормаль к поверхности & — \\
$\partial T/\partial n$ & производная по направлению нормали, $\Grad T\cdot\vec n$ & К/м \\
$\Grad T$ & градиент: вектор $(T_x,T_y,T_z)$, направлен в сторону наибольшего роста $T$ & К/м \\
$\Div\vec q$ & расходимость: $\partial q_x/\partial x+\partial q_y/\partial y+\partial q_z/\partial z$, мощность «источника» поля в точке & Вт/м$^3$ \\
$\Delta$ & оператор Лапласа, $\Delta T=\Div\Grad T=T_{xx}+T_{yy}+T_{zz}$ & 1/м$^2$ \\
$\Omega$ & область, занятая цилиндром (трёхмерная) & — \\
$\Pi$ & прямоугольник $0<r<a$, $0<z<h$ — половина осевого сечения & — \\
$\mu(z)$ & заданная температура боковой поверхности (как $\mu$ в условии I~рода из лекций) & К \\
$A$ & постоянная в примере неограниченного решения $A\ln(r/a)$ & К \\
$T_1$, $T_2$ & два решения: в доказательстве единственности — одной и той же задачи, в оценке устойчивости — задач с граничными функциями $\mu_1$, $\mu_2$ & К \\
$w$ & разность двух решений (в доказательстве единственности) & К \\
$a_{11}$, $a_{12}$, $a_{22}$ & коэффициенты при старших производных в общем виде уравнения; с радиусом $a$ не связаны & — \\
$D$ & $a_{12}^2-a_{11}a_{22}$, определяет тип уравнения & — \\
$\Phi$ & младшие члены уравнения (в лекциях — $F$); у нас $\Phi=\frac1r\,T_r$ & К/м$^2$ \\
$\lambda_1$, $\lambda_2$, $\lambda_3$ & переменные характеристической квадратичной формы (лекция~1) & — \\
$\xi$, $\zeta$ & безразмерные координаты $r/a$, $z/h$ & — \\
$\theta$ & безразмерная температура $T/T_0$ & — \\
$\mu_0(\zeta)$ & безразмерная температура боковой поверхности & — \\
$\varepsilon$ & малое изменение граничной функции (в оценке устойчивости) & К \\
$\alpha$ & произвольный угол поворота вокруг оси & рад \\
$u$, $f$ & искомая функция и свободный член в обозначениях лекций; в отчёте — только в формуле $u_t=a^2\Delta u+f$ (там $f=F/(C\rho)$) & — \\
\bottomrule
\end{longtable}

\textbf{Чем обозначения отчёта отличаются от лекций и почему.}

\begin{longtable}{@{}>{\raggedright\arraybackslash}p{4.6cm}>{\raggedright\arraybackslash}p{2.2cm}>{\raggedright\arraybackslash}p{2.6cm}>{\raggedright\arraybackslash}p{6.2cm}@{}}
\toprule
Величина & В лекциях & В отчёте & Причина \\
\midrule
\endhead
искомая функция & $u$ & $T$ & так в условии задачи \\
температуропроводность & $a^2$ & $\chi$ & $a$ занято радиусом \\
удельная теплоёмкость & $c$ & $C$ & $c$ занято высотой скачка \\
область & $D$ & $\Omega$, $\Pi$ & $D$ занято дискриминантом \\
радиус круга & $R$ & $a$ & так в условии задачи \\
полярный угол & $\theta$ & $\varphi$ & $\theta$ — безразмерная температура \\
младшие члены уравнения & $F$ & $\Phi$ & $F$ — источники тепла \\
\bottomrule
\end{longtable}

В вопросах ниже встречаются ещё буквы из лекций: $\Gamma$ — граница
области, $\overline D=D\cup\Gamma$ — замкнутая область (здесь $D$ — область,
а не дискриминант); $\alpha$, $\beta$, $\omega$, $\nu$ — коэффициенты и
правые части граничных условий
$(\alpha u+\beta\,\partial u/\partial n)|_\Gamma=\omega$ и
$\partial u/\partial n|_\Gamma=\nu$ (не путать с углом поворота $\alpha$).

% ============================================================
\section{Разбор формул и рисунков отчёта}
\label{sec:formulas}

\subsection{Раздел 1 «Постановка задачи»}

\fhead{\rf{eq:task}}{условие на боковой поверхности}
Запись $T|_{r=a}$ читается «значение $T$ при $r=a$», то есть на боковой
поверхности. Фигурная скобка задаёт кусочно-постоянную функцию: на нижней
части боковой поверхности $0<z<c$ (нижнем «поясе») температура нулевая, на
верхней части $c\le z<h$ (верхнем «поясе») равна $T_0$. Неравенства $0<z$
и $z<h$ строгие, потому что при $z=0$ и $z=h$ боковая поверхность
заканчивается рёбрами цилиндра (окружностями): они принадлежат сразу и
боковой поверхности, и торцу. Знак «$\le$» при $z=c$ на решение не влияет
(см. вопрос~\ref{q:pointc}).

\textbf{Про «ноль».} Это не абсолютный нуль, а начало отсчёта. Уравнение и
условие на торцах содержат только производные, поэтому к температуре можно
прибавить постоянную. Если нижний пояс имеет температуру $T_{\text{н}}$,
верхний — $T_{\text{в}}$, то $T-T_{\text{н}}$ решает нашу задачу с
$T_0=T_{\text{в}}-T_{\text{н}}$.

\subsection{Раздел 2.2 «Уравнение стационарной теплопроводности»}

\fhead{\rf{eq:fourier}}{закон Фурье}
$\vec q=-k\,\Grad T$. Опытный закон: плотность потока тепла пропорциональна
градиенту температуры. $\vec q$ показывает, куда и сколько тепла течёт:
его модуль — количество тепла через единичную площадку, перпендикулярную
потоку, за единицу времени. Градиент направлен в сторону роста температуры,
а тепло идёт от горячего к холодному, отсюда минус. $k>0$ — свойство
материала. То, что $k$ — число, а не матрица, и есть изотропность.

\fhead{\rf{eq:balance}}{баланс тепла для объёма $V$}
Закон сохранения энергии. Слева: в элементе $dV$ запасено тепло
$C\rho T\,dV$ (теплоёмкость $\times$ масса $\times$ температура), его
изменение в единицу времени — $C\rho\,T_t\,dV$; интеграл — для всего объёма.
Справа первое слагаемое: $\vec q\cdot\vec n\,dS$ — тепло, \emph{выходящее}
через площадку $dS$ (нормаль внешняя), поэтому приток — с минусом. Второе
слагаемое — тепло, выделяемое источниками. В лекции~2 тот же баланс записан
для стержня: $Q_3=Q_1+Q_2$.

\fhead{\rf{eq:ostrogradsky}}{формула Остроградского}
Поток вектора через замкнутую поверхность равен интегралу от расходимости
по объёму внутри. Она позволяет все три слагаемых баланса записать как
интегралы по одному и тому же объёму. Другое название — формула
Гаусса–Остроградского. Условия: поверхность $S$ замкнутая и
кусочно-гладкая, нормаль внешняя, вектор непрерывно дифференцируем в $V$
вплоть до поверхности. Эта же формула используется второй раз — в
доказательстве единственности~\rf{eq:energy}.

\fhead{\rf{eq:balance-vol}}{баланс одним объёмным интегралом}
Получается из~\rf{eq:balance}: поверхностный интеграл заменён объёмным по
формуле~\rf{eq:ostrogradsky}, вместо $\vec q$ подставлен закон
Фурье~\rf{eq:fourier} (минус на минус даёт плюс: $-\Div\vec q=\Div(k\,\Grad T)$),
все слагаемые перенесены в левую часть.

\textbf{Почему из равенства интеграла нулю следует равенство нулю
подынтегральной функции.} Объём $V$ любой. Если бы функция в какой-то точке
была, например, положительна, то по непрерывности она была бы положительна
и в малом шарике вокруг этой точки. Взяв этот шарик в качестве $V$, получили
бы положительный интеграл — противоречие.

\fhead{\rf{eq:heat-general}}{уравнение теплопроводности в общем виде}
$C\rho\,T_t=\Div(k\,\Grad T)+F$. Справедливо и для неоднородного тела, где
$k$ зависит от точки.

\fhead{\rf{eq:lap-cart}}{оператор Лапласа в декартовых координатах}
Сумма вторых производных по трём координатам. Получается так:
$\Grad T=(T_x,T_y,T_z)$, а расходимость этого вектора равна
$T_{xx}+T_{yy}+T_{zz}$, то есть $\Delta T=\Div\Grad T$. Оператор — это
$\Delta$ (правило), а $\Delta T$ — результат его применения к функции~$T$.

\fhead{\rf{eq:heat}}{уравнение теплопроводности для однородного тела}
При $k=\mathrm{const}$ множитель $k$ выносится:
$\Div(k\,\Grad T)=k\,\Div\Grad T=k\,\Delta T$. Делим на $C\rho$ и получаем
$T_t=\chi\,\Delta T+F/(C\rho)$. Это то же уравнение, что в лекции~2:
$u_t=a^2\Delta u+f$, только $a^2$ названо $\chi$, а $f=F/(C\rho)$.
$\chi=k/(C\rho)$ — коэффициент температуропроводности, м$^2$/с: показывает,
как быстро выравнивается температура.

\fhead{\rf{eq:laplace}}{уравнение Лапласа}
Стационарный режим: $T_t=0$. Источников нет: $F=0$. Остаётся
$\chi\,\Delta T=0$, и на $\chi\ne0$ можно сократить. Смысл: в каждый малый
объём входит столько же тепла, сколько выходит. Свойства материала $k$, $C$,
$\rho$ сократились — стационарное поле от них не зависит (от $k$ зависит
только величина потока $\vec q$). Функции, удовлетворяющие уравнению
Лапласа, называются гармоническими.

\subsection{Раздел 2.3 «Цилиндрические координаты и осевая симметрия»}

\fhead{\rf{eq:cyl}}{связь декартовых и цилиндрических координат}
$r$ и $\varphi$ — полярные координаты проекции точки на плоскость~$Oxy$,
$z$ — высота. Обратно: $r=\sqrt{x^2+y^2}$. Элемент объёма
$dV=r\,dr\,d\varphi\,dz$. Ось $Oz$ направлена от нижнего торца к верхнему —
от этого зависит знак нормали на торцах.

\fhead{\rf{eq:Omega}}{область $\Omega$}
Множество точек цилиндра: $r$ от 0 до $a$, угол любой, $z$ от 0 до $h$.
Неравенства $r<a$, $0<z<h$ строгие: область открытая, граница в неё не
входит (уравнение выполняется внутри, условия — на границе). $r\ge0$ с
равенством: ось принадлежит телу.

\fhead{\rf{eq:lap-cyl}}{оператор Лапласа в цилиндрических координатах}
Получается из $T_{xx}+T_{yy}+T_{zz}$ заменой переменных~\rf{eq:cyl} по
правилу дифференцирования сложной функции (вывод — в п.~\ref{der:lap}).
Слагаемое с $z$ не меняется, потому что $z$ не преобразуется. В учебнике
Смирнова (п.~131) формула записана в общем виде для ортогональных
координат. Коэффициенты оператора зависят от $r$, но не от $\varphi$ — это
используется в доказательстве симметрии.

\fhead{\rf{eq:axisym}}{осевая симметрия}
$T$ не зависит от $\varphi$. Простое объяснение: цилиндр симметричен
относительно оси, условия на поверхности от угла не зависят — ни одно
значение угла не выделено. Строгое (подраздел~2.6 отчёта): если
$T(r,\varphi,z)$ — решение трёхмерной задачи~\rf{eq:problem3d}, то
повёрнутая функция $T(r,\varphi+\alpha,z)$ — тоже её решение, потому что
коэффициенты оператора~\rf{eq:lap-cyl} и граничные условия от $\varphi$ не
зависят. Решение единственно, значит, обе функции совпадают при
любом~$\alpha$, то есть $T$ от $\varphi$ не зависит. Важны три вещи: тело
вращения, условия без~$\varphi$ и единственность, причём единственность
именно трёхмерной задачи — она и доказана в отчёте формулой~\rf{eq:energy}.

\fhead{\rf{eq:main-div} и \rf{eq:main}}{основное уравнение}
В~\rf{eq:lap-cyl} отбрасываем слагаемое с $\varphi$ (оно равно нулю).
Запись~\rf{eq:main-div} называется дивергентной; раскрывая производную
произведения $\frac{\partial}{\partial r}(r\,T_r)=T_r+r\,T_{rr}$ и деля
на~$r$, получаем~\rf{eq:main}. Обе записи — одно и то же уравнение. Первая
удобна для баланса тепла и для разностной схемы, вторая — для определения
типа и разделения переменных. Слагаемое $\frac1r T_r$ отличает
осесимметричную задачу от плоской: площадь цилиндрической поверхности
радиуса~$r$ растёт пропорционально~$r$ (вывод из баланса —
в п.~\ref{der:ring}).

\fhead{\rf{eq:Pi}}{прямоугольник $\Pi$}
Так как $T$ не зависит от $\varphi$, достаточно знать её в одной
полуплоскости $\varphi=\mathrm{const}$. Сечение цилиндра такой
полуплоскостью — прямоугольник: по горизонтали $r$ от 0 до $a$, по
вертикали $z$ от 0 до $h$. Задача из трёхмерной стала двумерной.

\subsection{Раздел 2.4 «Граничные условия»}

\fhead{\rf{eq:bc-side}}{условие на боковой поверхности}
Это условие~\rf{eq:task}, записанное через функцию $\mu(z)$. Задано
значение самой искомой функции на границе — условие I~рода (Дирихле),
в лекциях $u|_\Gamma=\mu$. Неоднородное, так как $\mu\not\equiv0$.
Физически: поверхность прижата к двум термостатам, которые держат её
температуру независимо от того, что происходит внутри. $\mu(z)$
кусочно-постоянна: интервал $0<z<h$ разбит на два куска, на каждом функция
постоянна. В точке $z=c$ разрыв первого рода: существуют конечные пределы
слева (0) и справа ($T_0$), но они не равны; их разность $T_0$ — величина
скачка. Разрыв второго рода — когда хотя бы один из односторонних пределов
бесконечен или не существует.

\fhead{\rf{eq:flux-n}}{поток через поверхность}
$q_n=\vec q\cdot\vec n$ — проекция вектора потока на внешнюю нормаль, то
есть тепло, выходящее через единицу площади за единицу времени. Подставляем
закон Фурье: $q_n=-k\,\Grad T\cdot\vec n$. Скалярное произведение градиента
на единичный вектор — производная по направлению этого вектора:
$\Grad T\cdot\vec n=\partial T/\partial n$.

\fhead{\rf{eq:bc-ends}}{условия на торцах}
Теплоизоляция: $q_n=0$, а $k\ne0$, значит, $\partial T/\partial n=0$.
На верхнем торце нормаль $\vec n=(0,0,1)$, $\partial T/\partial n=T_z$.
На нижнем $\vec n=(0,0,-1)$, $\partial T/\partial n=-T_z$. Условие
однородное, поэтому знак не важен: $T_z=0$ на обоих торцах. Задана
нормальная производная — условие II~рода (Неймана), в лекциях
$\partial u/\partial\vec n|_\Gamma=\nu$ с $\nu=0$. Геометрически: изотермы
подходят к торцам под прямым углом. Неравенство $r<a$ строгое: на ребре
$r=a$ нормаль не определена, и ни одно из двух условий там не задаётся
(условие на боковой поверхности поставлено при $0<z<h$). Температура ребра
получается по непрерывности с боковой поверхности: 0 на нижнем ребре, $T_0$
на верхнем.

\fhead{\rf{eq:bc-axis}}{ограниченность на оси}
Ось — не поверхность тела, на ней ничего не задано. Но в переменных $(r,z)$
отрезок $r=0$ — сторона прямоугольника~$\Pi$, а в уравнении~\rf{eq:main}
стоит $1/r$. У такого уравнения есть «лишние» решения, например $\ln r$:
$(\ln r)''+\frac1r(\ln r)'=-\frac{1}{r^2}+\frac{1}{r^2}=0$. В отчёте
записано $A\ln(r/a)$ — то же решение, умноженное на постоянную $A$ и
сдвинутое на постоянную, чтобы под логарифмом стояла безразмерная величина.
При $r\to0$ оно уходит в бесконечность и соответствует источнику или стоку
тепла на оси (подробно — п.~\ref{der:axis}), а их нет. Условие
ограниченности такие решения отсекает. В лекции~6 такое же условие
$|u(0,t)|<\infty$ ставится в центре круглой мембраны. Неравенства
$0\le z\le h$ нестрогие: ограниченность нужна на всей оси, включая центры
торцов; в отличие от рёбер, это обычные точки торцов, и условие $T_z=0$ в
них действует.

\fhead{\rf{eq:bc-axis-n}}{следствие: $T_r=0$ на оси}
Если идти по прямой через ось, то по одну сторону от оси температура
$T(r,z)$ и по другую тоже $T(r,z)$ — функция чётная относительно оси.
Производная гладкой чётной функции в нуле равна нулю. Это не дополнительное
условие, а следствие; им удобно пользоваться в разностной схеме.

\subsection{Раздел 2.5 «Математическая постановка задачи»}

\fhead{\rf{eq:problem3d}}{постановка в трёхмерной области}
Задача до использования симметрии: искомая функция $T(r,\varphi,z)$,
уравнение Лапласа с оператором~\rf{eq:lap-cyl}, условия на торцах, на
боковой поверхности и на оси. Последняя строка — условие периодичности
$T(r,\varphi+2\pi,z)=T(r,\varphi,z)$: углы $\varphi$ и $\varphi+2\pi$ задают
одну и ту же точку. В лекции~6 так же ставится задача для круга:
ограниченность в центре и периодичность по углу — последние строки системы.
Граничные условия от $\varphi$ не зависят, поэтому решение осесимметрично,
и остаётся задача~\rf{eq:problem}.

\fhead{\rf{eq:problem}}{итоговая постановка}
Четыре строки: уравнение в области; два условия II~рода на торцах; условие
I~рода на боковой поверхности; ограниченность на оси. Справа от каждой
строки — где она действует. Оформление как в лекциях: система под фигурной
скобкой, после неё «где … — искомая функция, … — заданные параметры».
Проверка числа условий: по $z$ уравнение второго порядка — два условия
($z=0$ и $z=h$); по $r$ тоже второго — два условия ($r=a$ и ограниченность
при $r=0$).

\textbf{Про класс функций.} В лекции~2: $u\in C^2(D)\cap C(\overline D)$.
$C^2(D)$ — внутри области есть непрерывные производные до второго порядка
(чтобы уравнение имело смысл). $C(\overline D)$ — функция непрерывна
в замкнутой области, включая границу (чтобы граничное условие понималось как
предел изнутри). У нас $\mu(z)$ разрывна в точке $z=c$; если бы решение было
непрерывно вплоть до границы, его граничные значения были бы непрерывны —
противоречие. Поэтому требуем меньше: ограниченность и непрерывность вплоть
до границы всюду, кроме точки $(a,c)$ (в пространстве это окружность).
Для условий II~рода на торцах нужна ещё непрерывность производной $T_z$
вплоть до торцов: условие $T_z=0$ понимается как предел $T_z$ при подходе
к торцу изнутри.

\textbf{Про гладкость границы.} В лекции~2 граница $\Gamma\in C^2$ —
гладкая. У цилиндра она кусочно-гладкая: на рёбрах нормаль не определена.
Поэтому условия поставлены на открытых частях границы ($r<a$ на торцах,
$0<z<h$ на боковой поверхности), а на рёбрах не задано ничего.

\subsection{Рисунки отчёта}

\textbf{Рисунок 1 — цилиндр и условия на его поверхности.} Это сам цилиндр
в пространстве. Точка $O$ — начало координат, центр нижнего торца. Ось $z$
идёт по оси цилиндра вверх, ось $r$ — по радиусу; штрихпунктирная линия —
ось цилиндра. Буква $a$ на оси $r$ — место, где ось $r$ выходит на боковую
поверхность, то есть $r=a$. Размерные линии слева: $h$ — высота цилиндра,
$c$ — высота окружности $z=c$, на которой температура боковой поверхности
меняется скачком. Штриховые дуги — невидимые (задние) половины окружностей
$z=0$ и $z=c$. Серая заливка — верхняя часть боковой поверхности
$c\le z<h$, где $T=T_0$; нижняя часть $0<z<c$ не закрашена, там $T=0$.
Наклонная штриховка верхнего торца — условное обозначение теплоизоляции.
Нижний торец тоже изолирован, но он обращён вниз и на рисунке не виден,
поэтому штриховки на нём нет. Условие $\partial T/\partial z=0$ подписано
выносками у обоих торцов.

\textbf{Рисунок 2 — расчётная область и граничные условия.} Это
прямоугольник $\Pi$ в координатах $(r,z)$ — половина осевого сечения; весь
цилиндр получается вращением прямоугольника вокруг оси $z$. Светло-серая
заливка — сама область $\Pi$. Внутри записано уравнение в
виде~\rf{eq:main-div}; в постановке~\rf{eq:problem} оно стоит в равносильном
виде~\rf{eq:main}. Стороны прямоугольника:
\begin{itemize}
  \item нижняя $z=0$ и верхняя $z=h$ — торцы. Штриховка \emph{снаружи} от
  них обозначает теплоизоляцию, рядом подписано условие
  $\partial T/\partial z=0$;
  \item правая $r=a$ — боковая поверхность. Ниже точки $z=c$ линия чёрная,
  там $T=0$; выше — утолщённая серая, там $T=T_0$ (на рисунке~1 серым
  закрашена та же часть поверхности);
  \item левая $r=0$ — ось цилиндра. Она штрихпунктирная, потому что это ось
  симметрии, а не поверхность тела: граничного условия на ней нет, стоит
  только требование ограниченности $|T|<\infty$.
\end{itemize}
Жирная точка на правой стороне — точка $(a,c)$, в которой граничная
температура меняется скачком (в пространстве это окружность $r=a$, $z=c$).
Штриховая горизонталь — уровень $z=c$; она проведена для наглядности: внутри
области на ней ничего не задано, решение там гладкое. Стрелки $\vec n$ —
внешние нормали к поверхности цилиндра: на верхнем торце вверх, на нижнем
вниз, на боковой поверхности — по радиусу от оси. На оси нормали нет, потому
что ось не является поверхностью. Буквы $O$, $a$, $c$, $h$ у осей — начало
координат и значения координат.

\subsection{Раздел 2.6 «Тип уравнения и свойства задачи»}

\fhead{\rf{eq:general}}{общий вид уравнения второго порядка}
Это формула из лекции~1 (там $a_{11}u_{xx}+2a_{12}u_{xy}+a_{22}u_{yy}+F=0$),
записанная для наших переменных: вместо $x$, $y$ стоят $r$, $z$, вместо
$u$ — $T$, вместо $F$ — $\Phi$. Слова «линейное относительно старших
производных» значат, что вторые производные входят линейно, а всё остальное
собрано в $\Phi$. Двойка при $a_{12}$: смешанные производные $T_{rz}$
и $T_{zr}$ равны и складываются. У нас $a_{11}=1$ (при $T_{rr}$), $a_{12}=0$
(смешанной производной нет), $a_{22}=1$ (при $T_{zz}$),
$\Phi=\frac1r\,T_r$.

\fhead{\rf{eq:discr}}{тип уравнения}
Правило из лекции~1: считаем $D=a_{12}^2-a_{11}a_{22}$; $D>0$ —
гиперболический тип, $D=0$ — параболический, $D<0$ — эллиптический. У нас
$D=0-1\cdot1=-1<0$. Младший член $\frac1r T_r$ входит в $\Phi$ и на тип не
влияет: тип определяют только старшие производные. $D$ — дискриминант
характеристического уравнения $a_{11}\,dy^2-2a_{12}\,dx\,dy+a_{22}\,dx^2=0$
(лекция~1); у нас $x=r$, $y=z$, получается $dz^2+dr^2=0$, то есть
$(dz/dr)^2=-1$ — вещественных решений нет, вещественных характеристик нет.
Для трёхмерного уравнения $\Delta T=0$ тип определяют по квадратичной форме
$\lambda_1^2+\lambda_2^2+\lambda_3^2$: все коэффициенты равны $+1$ — не
нули и одного знака, тип эллиптический.

\fhead{\rf{eq:energy}}{единственность}
Подробный вывод — в п.~\ref{der:uniq}. Идея: разность двух решений $w$
удовлетворяет той же задаче с нулевыми данными; интеграл от квадрата её
градиента по цилиндру равен интегралу по поверхности от
$w\,\partial w/\partial n$, а он нулевой; значит, $w$ постоянна, а так как
на боковой поверхности $w=0$, то $w\equiv0$. Доказательство проведено для
трёхмерной задачи~\rf{eq:problem3d} — без предположения о симметрии.

\fhead{\rf{eq:bounds}}{границы температуры}
Принцип максимума: гармоническая функция, отличная от постоянной, не может
иметь максимум или минимум внутри области — только на границе. Физически:
без источников внутри не может быть точки горячее всех соседних. Торцы
после чётного продолжения становятся внутренними точками, поэтому экстремумы
— на боковой поверхности, где $T$ равно 0 или $T_0$. Подробнее —
п.~\ref{der:max}.

\subsection{Раздел 2.7 «Безразмерная форма задачи»}

\fhead{\rf{eq:dimless-vars}}{безразмерные переменные}
Каждую величину делим на её характерный масштаб: радиус — на $a$, высоту —
на $h$, температуру — на $T_0$. Новые переменные меняются от 0 до~1.

\fhead{\rf{eq:dimless}}{безразмерная задача}
Вывод — в п.~\ref{der:dimless}. Главное: в задаче остались только два
числа, $a/h$ и $c/h$. Было четыре размерных параметра ($a$, $h$, $c$, $T_0$)
и две независимые размерности (длина и температура): $4-2=2$.

\fhead{\rf{eq:scaling}}{зависимость от $T_0$}
Задача линейна: если граничную температуру умножить на число, решение
умножится на то же число. Поэтому $T=T_0\,\theta$, и достаточно решить
задачу с единичным скачком.

% ============================================================
\section{Выводы, которые могут попросить воспроизвести}
\label{sec:derivations}

\subsection{Уравнение теплопроводности}
\label{der:heat}
\begin{enumerate}
  \item Закон Фурье: $\vec q=-k\,\Grad T$.
  \item Баланс для произвольного объёма $V$ с поверхностью $S$:
  \[
  \iiint\limits_V C\rho\,T_t\,dV=-\iint\limits_S\vec q\cdot\vec n\,dS+\iiint\limits_V F\,dV .
  \]
  \item Формула Остроградского:
  $\iint_S\vec q\cdot\vec n\,dS=\iiint_V\Div\vec q\,dV$. Подставляем
  $\vec q=-k\,\Grad T$:
  \[
  \iiint\limits_V\bigl(C\rho\,T_t-\Div(k\,\Grad T)-F\bigr)\,dV=0 .
  \]
  \item $V$ произволен $\Rightarrow$ $C\rho\,T_t=\Div(k\,\Grad T)+F$.
  \item $k=\mathrm{const}$ $\Rightarrow$ $T_t=\chi\,\Delta T+F/(C\rho)$,
  $\chi=k/(C\rho)$.
  \item $T_t=0$, $F=0$ $\Rightarrow$ $\Delta T=0$.
\end{enumerate}
В лекции~2 то же самое сделано для стержня: $Q_3=Q_1+Q_2$, где $Q_1$ —
тепло от потока, $Q_2$ — от источников, $Q_3$ — на нагрев.

\subsection{Оператор Лапласа в цилиндрических координатах}
\label{der:lap}

\textbf{Короткий вывод для осесимметричной функции} $T=T(r,z)$,
$r=\sqrt{x^2+y^2}$. Так как $\partial r/\partial x=x/r$, то
\[
T_x=T_r\,\frac{x}{r},\qquad
T_{xx}=T_{rr}\,\frac{x^2}{r^2}+T_r\Bigl(\frac{1}{r}-\frac{x^2}{r^3}\Bigr),
\]
и точно так же
$T_{yy}=T_{rr}\,\dfrac{y^2}{r^2}+T_r\Bigl(\dfrac{1}{r}-\dfrac{y^2}{r^3}\Bigr)$.
Складываем, учитывая $x^2+y^2=r^2$:
\[
T_{xx}+T_{yy}=T_{rr}+T_r\Bigl(\frac{2}{r}-\frac{1}{r}\Bigr)=T_{rr}+\frac{1}{r}\,T_r .
\]
Добавляем $T_{zz}$ — получаем левую часть~\rf{eq:main}.

\textbf{Полный вывод (с углом).} Из $r=\sqrt{x^2+y^2}$,
$\operatorname{tg}\varphi=y/x$ (дифференцируем эти равенства по $x$ и по
$y$):
\[
r_x=\cos\varphi,\quad r_y=\sin\varphi,\quad
\varphi_x=-\frac{\sin\varphi}{r},\quad \varphi_y=\frac{\cos\varphi}{r}.
\]
Поэтому
\[
\pd{}{x}=\cos\varphi\,\pd{}{r}-\frac{\sin\varphi}{r}\,\pd{}{\varphi},\qquad
\pd{}{y}=\sin\varphi\,\pd{}{r}+\frac{\cos\varphi}{r}\,\pd{}{\varphi}.
\]
Применяем каждый из операторов дважды:
\[
T_{xx}=\cos^2\!\varphi\,T_{rr}-\frac{2\sin\varphi\cos\varphi}{r}\,T_{r\varphi}
+\frac{\sin^2\!\varphi}{r}\,T_r+\frac{2\sin\varphi\cos\varphi}{r^2}\,T_\varphi
+\frac{\sin^2\!\varphi}{r^2}\,T_{\varphi\varphi},
\]
\[
T_{yy}=\sin^2\!\varphi\,T_{rr}+\frac{2\sin\varphi\cos\varphi}{r}\,T_{r\varphi}
+\frac{\cos^2\!\varphi}{r}\,T_r-\frac{2\sin\varphi\cos\varphi}{r^2}\,T_\varphi
+\frac{\cos^2\!\varphi}{r^2}\,T_{\varphi\varphi}.
\]
При сложении смешанные слагаемые сокращаются, а
$\sin^2\!\varphi+\cos^2\!\varphi=1$:
\[
T_{xx}+T_{yy}=T_{rr}+\frac1r\,T_r+\frac{1}{r^2}\,T_{\varphi\varphi}
=\frac1r\,\pd{}{r}\Bigl(r\,\pd{T}{r}\Bigr)+\frac{1}{r^2}\,\pdd{T}{\varphi}.
\]

\subsection{Основное уравнение из баланса тепла в кольце}
\label{der:ring}
Этот вывод не использует декартовы координаты и объясняет множитель $1/r$.
Берём кольцо: радиус от $r$ до $r+dr$, высота от $z$ до $z+dz$, угол любой.
\begin{itemize}
  \item Через внутреннюю цилиндрическую стенку (площадь $2\pi r\,dz$)
  в направлении роста $r$ проходит тепло $Q_r(r)=-k\,T_r\cdot2\pi r\,dz$.
  Оно входит в кольцо. Через наружную стенку выходит $Q_r(r+dr)$. Приток:
  \[
  Q_r(r)-Q_r(r+dr)=2\pi k\,dz\Bigl[(r\,T_r)\big|_{r+dr}-(r\,T_r)\big|_{r}\Bigr]
  =2\pi k\,\pd{}{r}\Bigl(r\,\pd{T}{r}\Bigr)dr\,dz .
  \]
  \item Через нижнюю и верхнюю грани (площадь каждой $2\pi r\,dr$) приток
  равен $2\pi k\,r\,T_{zz}\,dr\,dz$.
  \item Режим стационарный, источников нет: сумма притоков равна нулю,
  \[
  \pd{}{r}\Bigl(r\,\pd{T}{r}\Bigr)+r\,\pdd{T}{z}=0 .
  \]
  Делим на $r$ — получаем~\rf{eq:main-div}.
\end{itemize}
Вывод: в баланс входит не плотность радиального потока $-k\,T_r$, а полный
поток через цилиндрическую стенку — плотность, умноженная на площадь стенки
$2\pi r\,dz$, то есть величина, пропорциональная $r\,T_r$. Поэтому
в уравнении стоит $\frac{\partial}{\partial r}(r\,T_r)$, а не $T_{rr}$.
Сама величина $r\,T_r$ постоянна по $r$ только при чисто радиальном течении
тепла ($T_{zz}=0$); в нашей задаче её изменение компенсируется осевым
потоком: $(r\,T_r)_r=-r\,T_{zz}$.

\subsection{Условие теплоизоляции}
\label{der:ins}
Тепло, выходящее через площадку $dS$ за единицу времени:
$\vec q\cdot\vec n\,dS=-k\,\dfrac{\partial T}{\partial n}\,dS$.
Изоляция: это тепло равно нулю для любой площадки торца, значит,
$\partial T/\partial n=0$ во всех точках торца. Если бы через торец был
задан поток $q_0$, \emph{входящий} в тело, условие имело бы вид
$k\,T_z=q_0$ при $z=h$ и $-k\,T_z=q_0$ при $z=0$ — здесь знак нормали уже
важен.

\subsection{Условие на оси}
\label{der:axis}
Ищем решения уравнения~\rf{eq:main-div}, не зависящие от $z$:
$\frac1r(r\,T')'=0$ $\Rightarrow$ $r\,T'=C_1$ $\Rightarrow$
$T=C_1\ln r+C_2$. При $C_1\ne0$ температура на оси бесконечна. Через цилиндр
малого радиуса $r_0$ вокруг оси на единицу длины наружу проходит тепло
$-k\,T_r\cdot2\pi r_0=-2\pi k\,C_1$ — оно не стремится к нулю при
$r_0\to0$, то есть на оси находится линейный источник (при $C_1<0$) или сток
(при $C_1>0$) мощностью $2\pi k|C_1|$ на единицу длины. Например, для
$T=\ln r$ производная $T_r=1/r>0$, поток направлен к оси — это сток. Ни
источников, ни стоков нет, поэтому $C_1=0$. Требование $|T|<\infty$ делает
это автоматически. На этапе~2 по той же причине из двух решений радиального
уравнения $I_0$ и $K_0$ останется только $I_0$ ($K_0$ — функция
Макдональда, второе решение модифицированного уравнения Бесселя; при
$r\to0$ она ведёт себя как $-\ln r$. Не путать с $K_0$ у Смирнова, где так
обозначено второе решение обычного уравнения Бесселя; в лекции~6 оно
названо $Y_m$).

\subsection{Единственность решения}
\label{der:uniq}
Доказываем для трёхмерной задачи~\rf{eq:problem3d}. Пусть $T_1$, $T_2$ —
два её решения, $w=T_1-T_2$. Задача линейна, поэтому
\[
\Delta w=0\ \text{в }\Omega,\qquad w=0\ \text{на боковой поверхности},\qquad
\pd{w}{n}=0\ \text{на торцах}.
\]
\begin{enumerate}
  \item Производная произведения: $\Div(w\,\Grad w)=\Grad w\cdot\Grad w
  +w\,\Div\Grad w=|\Grad w|^2+w\,\Delta w$. Так как $\Delta w=0$, остаётся
  $\Div(w\,\Grad w)=|\Grad w|^2$.
  \item Формула Остроградского~\rf{eq:ostrogradsky} для вектора
  $w\,\Grad w$ и объёма $\Omega$: интеграл от расходимости равен потоку
  через поверхность цилиндра $S_\Omega$. Поток вектора $w\,\Grad w$ — это
  интеграл от $w\,\Grad w\cdot\vec n=w\,\partial w/\partial n$:
  \[
  \iiint\limits_{\Omega}|\Grad w|^{2}\,dV
  =\iint\limits_{S_\Omega} w\,\pd{w}{n}\,dS .
  \]
  \item Справа нуль: на боковой поверхности $w=0$, на торцах
  $\partial w/\partial n=0$.
  \item Слева под интегралом неотрицательная функция, а интеграл равен
  нулю, значит, $\Grad w=0$ всюду, $w$ постоянна. На боковой поверхности
  она равна нулю, значит, $w\equiv0$, то есть $T_1\equiv T_2$.
\end{enumerate}
Двумерная задача~\rf{eq:problem} — частный случай: её решение — это решение
трёхмерной задачи, не зависящее от $\varphi$. Поэтому и оно единственно.

\textbf{То же в координатах $(r,z)$.} Если попросят расписать без векторных
обозначений: $w=w(r,z)$,
$\frac{\partial}{\partial r}(r\,w_r)+r\,w_{zz}=0$. Умножаем на $w$,
интегрируем по $\Pi$ по частям:
\[
\int\limits_0^a w\,\pd{}{r}\Bigl(r\,\pd{w}{r}\Bigr)dr
=\Bigl[r\,w\,w_r\Bigr]_{r=0}^{r=a}-\int\limits_0^a r\,w_r^2\,dr,
\qquad
\int\limits_0^h w\,w_{zz}\,dz=\Bigl[w\,w_z\Bigr]_{z=0}^{z=h}-\int\limits_0^h w_z^2\,dz .
\]
Внеинтегральные слагаемые равны нулю: при $r=a$ — потому что $w=0$; при
$r=0$ — из-за множителя $r$; при $z=0$ и $z=h$ — потому что $w_z=0$.
Остаётся $\int_0^h\!\int_0^a(w_r^2+w_z^2)\,r\,dr\,dz=0$. Это та же
формула~\rf{eq:energy}, делённая на $2\pi$: $dV=r\,dr\,d\varphi\,dz$,
$|\Grad w|^2=w_r^2+w_z^2$.

\textbf{Замечания.} 1) Условие I~рода на боковой поверхности убирает
произвол в постоянной: в чистой задаче Неймана решение определено с
точностью до константы, у нас — нет.

2) Оговорка из отчёта. Формула Остроградского требует, чтобы производные
$w$ были ограничены вплоть до поверхности. У $w$ граничные данные нулевые,
но сама $w$ — разность двух решений с разрывной $\mu$, а о каждом из них
около окружности $r=a$, $z=c$ известно лишь то, что оно ограничено
(градиент самого решения там растёт как $T_0/(\pi d)$, где $d$ — расстояние
до окружности). Поэтому заранее не известно, ограничены ли производные $w$,
и выкладка формально не обоснована. Строго единственность для разрывных
данных доказывают через принцип максимума — среди ограниченных решений
(Тихонов, Самарский, гл.~IV). На вопрос «докажите единственность» отвечать
выкладкой выше и добавлять эту оговорку; на вопрос «при чём тут разрыв,
если у $w$ данные нулевые?» — первыми тремя предложениями этого замечания.

\subsection{Принцип максимума и оценка $0\le T\le T_0$}
\label{der:max}
\begin{enumerate}
  \item Чётное продолжение. Положим $T(r,-z)=T(r,z)$. Производная по $z$
  входит в уравнение только в виде $T_{zz}$, поэтому оно не меняется при
  замене $z\to-z$; из-за условия $T_z=0$ при $z=0$ функция и её производная
  склеиваются непрерывно. Продолженная функция гармонична в цилиндре
  $-h<z<h$, и точки нижнего торца стали внутренними. То же — с верхним
  торцом.
  \item Принцип максимума (Смирнов, п.~204): гармоническая функция, не
  равная постоянной, не достигает наибольшего и наименьшего значений во
  внутренних точках.
  \item Пусть наибольшее значение достигается в точке торца $z=0$, $r<a$.
  Продолженная функция принимает те же значения, что и исходная, а эта
  точка для неё внутренняя. По принципу максимума функция тогда постоянна,
  что невозможно при $T_0\ne0$. То же — для торца $z=h$ и для наименьшего
  значения. Внутри цилиндра экстремума тоже нет. Остаётся боковая
  поверхность, где температура равна 0 или $T_0$. Отсюда при $T_0>0$
  получаем $0\le T\le T_0$, а внутри строго $0<T<T_0$. Торцы продолженного
  цилиндра не мешают: его граница в рассуждении не используется, нужно
  только, чтобы точка торца исходного цилиндра стала внутренней.
\end{enumerate}
Устойчивость: разность двух решений с граничными функциями $\mu_1$, $\mu_2$
— решение с граничной функцией $\mu_1-\mu_2$, поэтому
$|T_1-T_2|\le\max|\mu_1-\mu_2|$.

Чётное продолжение объясняет и будущий вид решения: продолженная функция
чётна и периодична по $z$ с периодом $2h$, поэтому раскладывается в ряд по
$\cos(\pi nz/h)$.

\subsection{Безразмерная форма}
\label{der:dimless}
Подставляем $r=a\xi$, $z=h\zeta$, $T=T_0\theta$:
\[
\pd{T}{r}=\frac{T_0}{a}\,\pd{\theta}{\xi},\qquad
\frac1r\,\pd{}{r}\Bigl(r\,\pd{T}{r}\Bigr)
=\frac{T_0}{a^2}\,\frac1\xi\,\pd{}{\xi}\Bigl(\xi\,\pd{\theta}{\xi}\Bigr),\qquad
\pdd{T}{z}=\frac{T_0}{h^2}\,\pdd{\theta}{\zeta}.
\]
Складываем и умножаем на $a^2/T_0$:
\[
\frac1\xi\,\pd{}{\xi}\Bigl(\xi\,\pd{\theta}{\xi}\Bigr)
+\Bigl(\frac{a}{h}\Bigr)^2\pdd{\theta}{\zeta}=0 .
\]
Условия: $T_z=0$ переходит в $\theta_\zeta=0$ при $\zeta=0$ и $\zeta=1$;
$T(a,z)=\mu(z)$ — в $\theta(1,\zeta)=\mu(h\zeta)/T_0=\mu_0(\zeta)$; линия
скачка $z=c$ — в $\zeta=c/h$.

\subsection{Что можно сказать о решении, не решая задачу}
\label{der:checks}
Эти факты пригодятся на этапах 2–4 для проверки формул и графиков.
\begin{itemize}
  \item Предельные случаи: при $c\to0$ вся боковая поверхность имеет
  температуру $T_0$ и $T\equiv T_0$; при $c\to h$ — $T\equiv0$.
  \item $0\le T\le T_0$ (при $T_0>0$).
  \item Средняя по высоте температура не зависит от $r$. Обозначим
  $m(r)=\frac1h\int_0^h T\,dz$ и проинтегрируем~\rf{eq:main-div} по $z$ от
  0 до $h$: интеграл от $T_{zz}$ равен $T_z|_0^h=0$, остаётся
  $(r\,m')'=0$. Ограниченное решение — постоянная. При $r=a$ она равна
  $\frac1h\int_0^h\mu\,dz=T_0(1-c/h)$. Физически: суммарный поток тепла
  через цилиндрическую поверхность любого радиуса на всю высоту равен нулю.
  \item Симметрия: $T(r,z;c)+T(r,h-z;h-c)=T_0$. При $c=h/2$ в среднем
  сечении $T(r,h/2)=T_0/2$.
  \item Все изотермы с уровнями от 0 до $T_0$ сходятся в точку $(a,c)$;
  к торцам и к оси подходят под прямым углом.
\end{itemize}
Справочно — к чему приведёт этап~2 (метод Фурье, $T=R(r)Z(z)$):
\[
T(r,z)=T_0\Bigl(1-\frac{c}{h}\Bigr)-\frac{2T_0}{\pi}\sum_{n=1}^{\infty}
\frac1n\,\sin\frac{\pi nc}{h}\,
\frac{I_0(\pi nr/h)}{I_0(\pi na/h)}\,\cos\frac{\pi nz}{h},
\]
где $I_0$ — модифицированная функция Бесселя нулевого порядка. По $z$ стоят
косинусы, потому что на торцах условия $T_z=0$; по $r$ — $I_0$, потому что
она ограничена на оси.

% ============================================================
\section{Вопросы и ответы}
\label{sec:qa}

\subsection{По условию задачи}

\q{Что значит «стационарное распределение»?}
Установившееся: температура в каждой точке не меняется со временем,
$\partial T/\partial t=0$. Такой режим наступает через большое время после
того, как на границе установили постоянные условия. Поэтому в постановке нет
ни времени, ни начального условия.

\q{Почему нет начального условия? Разве ответ не зависит от того, каким
цилиндр был вначале?}
Не зависит. В нестационарной задаче отклонение от стационарного
распределения затухает со временем по экспоненте, и при $t\to\infty$ от
начального распределения ничего не остаётся. Стационарное поле определяется
только условиями на границе.

\q{Что такое $T_0$ и что такое температура «ноль» на нижней части?}
$T_0$ — заданная постоянная. Ноль — начало отсчёта, а не абсолютный нуль:
к решению можно прибавить константу, поэтому температуру отсчитываем от
температуры нижнего пояса. $T_0$ — разность температур двух поясов.

\q{Что означают $a$, $h$, $c$? Какие на них ограничения?}
Радиус, высота, высота линии скачка. Все — длины, $a>0$, $h>0$, $0<c<h$.
При $c=0$ или $c=h$ скачка нет и решение — константа.

\q{От каких свойств материала зависит ответ?}
Ни от каких. $C\rho$ стоит при $T_t$, а $T_t=0$; $k$ в уравнении
сокращается, в условии изоляции $k\,T_z=0$ тоже. Нужно лишь, чтобы материал
был однородным и изотропным. От $k$ зависит только величина потока тепла.

\q{Опишите качественно, как течёт тепло.}
При $T_0>0$ тепло входит через верхний пояс боковой поверхности, проходит
через тело и выходит через нижний пояс. Через торцы не проходит. Сколько
входит, столько и выходит (поток через любую замкнутую поверхность внутри
тела равен нулю), иначе режим не был бы стационарным. Если спросят «сколько
именно»: из-за идеального скачка плотность потока на боковой поверхности
у линии $z=c$ растёт как $kT_0/(\pi d)$, где $d=|z-c|$, и полная мощность
через каждый пояс расходится логарифмически; при конечной ширине переходной
зоны она конечна. Основная часть тепла перетекает из верхнего пояса
в нижний вблизи линии скачка.

\q{Какие допущения приняты и где каждое использовано?}
Изотропность — закон Фурье~\rf{eq:fourier} с одним числом $k$. Однородность
и независимость $k$ от температуры — чтобы вынести $k$ из-под $\Div$ и
получить $\Delta T$. Нет источников — $F=0$ (сток — это источник с $F<0$,
поглощение тепла), уравнение Лапласа, а не Пуассона. Только
теплопроводность — применим закон Фурье. Стационарность — $T_t=0$.
Идеальная изоляция и идеальный термостат — условия II и I~рода в чистом
виде, а не теплообмен (III~род).

\q{Чем однородный материал отличается от изотропного?}
Однородный — свойства одинаковы во всех точках: $k$, $C$, $\rho$ не зависят
от координат. Изотропный — свойства одинаковы по всем направлениям:
теплопроводность задаётся одним числом $k$, и поток направлен точно против
градиента. Изотропность использована в законе Фурье~\rf{eq:fourier},
однородность — при переходе от~\rf{eq:heat-general} к~\rf{eq:heat}. В
анизотропном теле (кристалл, дерево, слоистый композит) вместо числа $k$
стоит матрица, и поток не параллелен градиенту.

\q{Зачем оговорено, что $k$ не зависит от температуры?}
Иначе уравнение $\Div\bigl(k(T)\,\Grad T\bigr)=0$ нелинейно: $k$ нельзя
вынести и сократить, перестают работать суперпозиция, разделение переменных
и пропорциональность решения величине $T_0$. При небольших перепадах
температуры $k$ меняется слабо, поэтому допущение оправдано.

\q{Какие ещё бывают способы переноса тепла и почему они не учтены?}
Конвекция (тепло переносит движущаяся среда) и излучение. Цилиндр — твёрдое
непрозрачное тело: вещество в нём не движется, излучение внутри не
распространяется. Остаётся теплопроводность, которую и описывает закон
Фурье.

\subsection{По выводу уравнения}

\q{Сформулируйте закон Фурье. Почему минус?}
$\vec q=-k\,\Grad T$: плотность потока тепла пропорциональна градиенту
температуры. Градиент смотрит в сторону роста температуры, тепло идёт
в сторону убывания.

\q{Что стоит в левой и правой частях баланса~\rf{eq:balance}?}
Слева — скорость изменения запаса тепла в объёме. Справа — приток тепла
через поверхность (минус, потому что $\vec q\cdot\vec n$ — это отток) и
тепло источников.

\q{Что такое формула Остроградского и зачем она здесь?}
Поток вектора через замкнутую поверхность равен интегралу от его
расходимости по объёму. Она переводит поверхностный интеграл в объёмный,
после чего всё равенство — один интеграл по $V$.

\q{Почему подынтегральная функция равна нулю?}
Интеграл равен нулю для любого объёма. Будь функция где-то ненулевой, она
сохраняла бы знак в малом шарике, и интеграл по этому шарику был бы
ненулевым.

\q{Что такое оператор Лапласа и каков его смысл?}
$\Delta T=\Div\Grad T=T_{xx}+T_{yy}+T_{zz}$. С точностью до множителя $k$
это приток тепла в единицу объёма. $\Delta T=0$ — в каждой точке приток
равен оттоку.

\q{Как из уравнения теплопроводности получилось уравнение Лапласа?}
$T_t=0$ (стационарность) и $F=0$ (нет источников). При наличии источников
получилось бы уравнение Пуассона $\Delta T=-F/k$.

\q{В лекциях $a^2=k/(c\rho)$, а у вас $a$ — радиус. Как так?}
Буквы $a$ и $c$ заняты условием задачи (радиус и высота скачка), поэтому
температуропроводность обозначена $\chi$ (она соответствует лекционному
$a^2$), теплоёмкость — $C$. В итоговую постановку они не входят.

\q{Проверьте размерности в уравнении~\rf{eq:heat-general}. Почему у $\chi$
размерность м$^2$/с?}
Каждое слагаемое — мощность в единице объёма, Вт/м$^3$:
\[
\Bigl[C\rho\,\pd{T}{t}\Bigr]=\frac{\text{Дж}}{\text{кг}\cdot\text{К}}\cdot
\frac{\text{кг}}{\text{м}^3}\cdot\frac{\text{К}}{\text{с}}
=\frac{\text{Вт}}{\text{м}^3},\qquad
\bigl[\Div(k\,\Grad T)\bigr]=\frac{1}{\text{м}}\cdot
\frac{\text{Вт}}{\text{м}\cdot\text{К}}\cdot\frac{\text{К}}{\text{м}}
=\frac{\text{Вт}}{\text{м}^3},
\]
и $F$ задано в Вт/м$^3$ (квадратные скобки означают «размерность величины»).
Для температуропроводности
\[
[\chi]=\Bigl[\frac{k}{C\rho}\Bigr]=\frac{\text{Вт}}{\text{м}\cdot\text{К}}:
\frac{\text{Дж}}{\text{м}^3\cdot\text{К}}=\frac{\text{м}^2}{\text{с}} ,
\]
так как Дж${}={}$Вт$\cdot$с. То же видно из~\rf{eq:heat}: $T_t$ измеряется
в К/с, $\Delta T$ — в К/м$^2$, поэтому их отношение $\chi$ — в м$^2$/с.

\q{Что такое гармоническая функция? Какие свойства знаете?}
Функция с непрерывными вторыми производными, удовлетворяющая уравнению
Лапласа. Свойства: значение в центре шара равно среднему по сфере; максимум
и минимум достигаются на границе области.

\subsection{По координатам и симметрии}

\q{Почему выбраны цилиндрические координаты?}
В них поверхность цилиндра — координатные поверхности $r=a$, $z=0$, $z=h$,
и граничные условия записываются просто. К тому же появляется возможность
разделить переменные.

\q{Где начало координат и куда направлена ось $z$?}
Начало — в центре нижнего торца, ось $z$ — по оси цилиндра вверх. Тогда
торцы — это $z=0$ и $z=h$, как в условии.

\q{Запишите оператор Лапласа в цилиндрических координатах. Как его
получить?}
Формула~\rf{eq:lap-cyl}. Получается заменой переменных по правилу
дифференцирования сложной функции; для осесимметричной функции вывод
занимает три строки (п.~\ref{der:lap}).

\q{Откуда слагаемое $\frac1r\,T_r$?}
Из геометрии: площадь цилиндрической поверхности $2\pi r\,dz$ растёт
с радиусом. Формально — из раскрытия $\frac1r(r\,T_r)_r$. В плоской задаче
его нет.

\q{Почему температура не зависит от $\varphi$? Докажите.}
Повёрнутая функция $T(r,\varphi+\alpha,z)$ удовлетворяет тому же уравнению
и тем же условиям: коэффициенты оператора~\rf{eq:lap-cyl} и граничные
условия от $\varphi$ не зависят. Решение трёхмерной
задачи~\rf{eq:problem3d} единственно — это доказано
формулой~\rf{eq:energy} без предположения о симметрии. Значит, повёрнутая
функция совпадает с исходной при любом $\alpha$, то есть $T$ от $\varphi$ не
зависит.

\q{Для какой задачи доказана единственность — двумерной или трёхмерной? Нет
ли круга: симметрия из единственности, а единственность — для уже
симметричной задачи?}
Круга нет. Единственность доказана для трёхмерной
задачи~\rf{eq:problem3d}, где зависимость от $\varphi$ не исключена:
формула~\rf{eq:energy} записана для объёма цилиндра $\Omega$ и его
поверхности. Порядок рассуждения: единственность в $\Omega$ $\to$ осевая
симметрия $\to$ двумерная задача~\rf{eq:problem}. Единственность двумерной
задачи — следствие: её решение есть решение трёхмерной, не зависящее от
$\varphi$.

\q{А если бы $T_0$ зависело от угла?}
Симметрия пропала бы: $T=T(r,\varphi,z)$, в уравнении осталось бы
$\frac{1}{r^2}T_{\varphi\varphi}$, условие периодичности по $\varphi$ уже
не выполнялось бы автоматически.

\q{В лекции~6 в задаче для круга есть условие периодичности по углу. Где
оно у вас?}
В трёхмерной постановке~\rf{eq:problem3d} — последняя строка: углам
$\varphi$ и $\varphi+2\pi$ отвечает одна и та же точка, поэтому
$T(r,\varphi+2\pi,z)=T(r,\varphi,z)$. После того как доказано, что $T$ от
$\varphi$ не зависит, оно выполняется автоматически, и в двумерную
постановку~\rf{eq:problem} его включать незачем.

\q{Почему уравнение записано при $0<r<a$, а не $0\le r<a$?}
При $r=0$ коэффициент $1/r$ не определён. Это особенность координат, а не
задачи: в декартовых координатах уравнение $T_{xx}+T_{yy}+T_{zz}=0$ на оси
ничем не выделено.

\q{Что такое $\Pi$ и почему задача стала двумерной?}
Прямоугольник $0<r<a$, $0<z<h$ — половина осевого сечения. Раз $T$ не
зависит от $\varphi$, достаточно найти её в одном таком сечении.

\subsection{По граничным условиям}

\q{Как фраза «поддерживается при температуре» превращается в формулу?}
Температура поверхности задана: $T(a,z)=\mu(z)$. Это условие I~рода
(Дирихле), неоднородное.

\q{Как фраза «теплоизолированные торцы» превращается в формулу?}
Тепло через торец не проходит: $q_n=-k\,\partial T/\partial n=0$, значит,
$\partial T/\partial z=0$ при $z=0$ и $z=h$. Условие II~рода (Неймана),
однородное.

\q{Что такое $\partial T/\partial n$? Куда направлена нормаль?}
Производная по направлению внешней нормали: $\Grad T\cdot\vec n$. На
верхнем торце нормаль вверх, $\partial T/\partial n=T_z$; на нижнем — вниз,
$\partial T/\partial n=-T_z$; на боковой поверхности — по радиусу,
$\partial T/\partial n=T_r$.

\q{На нижнем торце нормаль направлена вниз. Вы не потеряли знак?}
Там $\partial T/\partial n=-T_z$, но условие однородное: $-T_z=0$ и $T_z=0$
— одно и то же. При ненулевом заданном потоке знак был бы важен.

\q{Чем физически отличаются условия I, II и III~рода?}
I — задана температура поверхности. II — задан поток тепла через
поверхность (ноль при изоляции). III — теплообмен со средой по закону
Ньютона: поток пропорционален разности температур поверхности и среды,
$(\alpha u+\beta\,\partial u/\partial n)|_\Gamma=\omega$.

\q{Запишите ваши условия в общем виде
$(\alpha u+\beta\,\partial u/\partial n)|_\Gamma=\omega$ из лекций.}
На боковой поверхности $\alpha=1$, $\beta=0$, $\omega=\mu(z)$. На торцах
$\alpha=0$, $\beta=1$, $\omega=0$.

\q{Зачем условие на оси, если ось — не граница тела?}
У уравнения в цилиндрических координатах есть решения, неограниченные на
оси ($\ln r$). Они отвечают источнику или стоку тепла на оси, которых нет.
Условие ограниченности их отбрасывает. Это не граничное условие I–III рода:
на оси ничего не задаётся.

\q{Покажите, что на оси $T_r=0$.}
Вдоль прямой, проходящей через ось, осесимметричная функция чётна, а
производная гладкой чётной функции в нуле равна нулю.

\q{Можно ли в постановке вместо $|T|<\infty$ написать $T_r=0$ при $r=0$?}
Можно: оба условия отбрасывают одно и то же решение $C_1\ln r$ (у него
$T_r=C_1/r$ не стремится к нулю). В лекциях на оси ставится именно условие
ограниченности, поэтому в постановке оставлено оно, а $T_r=0$ используется
как его следствие в разностной схеме.

\q{Сколько условий нужно и почему именно столько?}
Для эллиптического уравнения второго порядка — одно условие в каждой точке
границы. У нас: на боковой поверхности — Дирихле, на торцах — Неймана.

\q{Что задано на рёбрах $r=a$, $z=0$ и $r=a$, $z=h$?}
Ребро принадлежит и торцу, и боковой поверхности, нормаль на нём не
определена. Поэтому условия заданы на открытых частях границы ($0<z<h$;
$r<a$), а на самом ребре — ничего. Противоречия нет: около рёбер $\mu(z)$
постоянна, её производная по $z$ равна нулю, что согласуется с $T_z=0$.
Температура на рёбрах получается по непрерывности: 0 и $T_0$.

\q{Важно ли, чему равна $\mu$ в самой точке $z=c$?\label{q:pointc}}
Нет. Изменение граничной функции в одной точке решение не меняет: решение
выражается через интегралы от $\mu$. В самой точке $(a,c)$ у решения нет
единого предела: предельное значение зависит от направления подхода; ряд
Фурье даст там $T_0/2$.

\q{А что, если задать на торцах ещё и температуру?}
Задача станет переопределённой. Наша задача уже имеет единственное решение;
его значения на торцах получаются сами и с произвольно заданной температурой
торцов, вообще говоря, не совпадут — решения не будет. Задавать на одной
части границы и $T$, и $\partial T/\partial n$, ничего не задавая на
остальной, — это задача Коши для уравнения Лапласа. Она некорректна по
другой причине: решение неустойчиво (пример Адамара, см. ниже).

\subsection{По постановке, типу и корректности}

\q{Сформулируйте задачу целиком.}
Найти функцию $T(r,z)$, дважды непрерывно дифференцируемую и ограниченную
в прямоугольнике $\Pi$: $0<r<a$, $0<z<h$, которая удовлетворяет в $\Pi$
уравнению $T_{rr}+\frac1r T_r+T_{zz}=0$, на торцах — условиям $T_z=0$ при
$z=0$ и $z=h$ ($0\le r<a$), на боковой поверхности — условию
$T(a,z)=\mu(z)$ при $0<z<h$, где $\mu=0$ при $0<z<c$ и $\mu=T_0$ при
$c\le z<h$, и на оси — условию ограниченности $|T(0,z)|<\infty$. Порядок тот
же, что у четырёх строк~\rf{eq:problem}.

\q{Запишите постановку в исходной трёхмерной области, до использования
симметрии.}
Это задача~\rf{eq:problem3d}. Найти $T(r,\varphi,z)$ в области $\Omega$:
$\Delta T=0$ с оператором~\rf{eq:lap-cyl}; $T_z=0$ при $z=0$ и $z=h$;
$T|_{r=a}=\mu(z)$; ограниченность при $r=0$; условие периодичности
$T(r,\varphi+2\pi,z)=T(r,\varphi,z)$.

\q{Определите тип уравнения.}
$a_{11}=1$, $a_{12}=0$, $a_{22}=1$, $D=a_{12}^2-a_{11}a_{22}=-1<0$ —
эллиптический тип во всей области. Младший член $\frac1r T_r$ на тип не
влияет.

\q{Что такое $D$? Запишите характеристическое уравнение. Каков канонический
вид вашего уравнения?}
$D=a_{12}^2-a_{11}a_{22}$ — дискриминант характеристического уравнения
$a_{11}\,dy^2-2a_{12}\,dx\,dy+a_{22}\,dx^2=0$ (лекция~1). У нас $x=r$,
$y=z$: $dz^2+dr^2=0$, $dz/dr=\pm i$ — вещественных характеристик нет, есть
два комплексно-сопряжённых интеграла $z\pm ir=\mathrm{const}$. Новые
переменные — действительная и мнимая части: $\xi=z$, $\eta=r$, то есть сами
$z$ и $r$. Значит, уравнение~\rf{eq:main} уже записано в каноническом виде
эллиптического типа $u_{\xi\xi}+u_{\eta\eta}+\Phi=0$, где $\Phi$ — младшие
члены. Так же выглядит канонический вид уравнения Трикоми в области
эллиптичности из лекции~1.

\q{А какого типа исходное уравнение $\Delta T=0$ с тремя переменными? Там
$D$ не посчитать.}
Для $n$ переменных тип определяют по характеристической квадратичной форме
$Q=\sum a_{ij}\lambda_i\lambda_j$ (лекция~1): её приводят к сумме квадратов
$\sum\beta_i\mu_i^2$ и смотрят на знаки $\beta_i$. Для уравнения Лапласа
$Q=\lambda_1^2+\lambda_2^2+\lambda_3^2$: все $\beta_i=1$ — не равны нулю и
одного знака, тип эллиптический. В лекции~2 так и записано: $\Delta u=f$ —
эллиптический тип. Для сравнения: $u_{tt}=a^2\Delta u$ — гиперболический
($\beta_i$ не равны нулю, но разных знаков), $u_t=a^2\Delta u$ —
параболический (есть $\beta_i=0$).

\q{Уравнение линейное? Однородное?}
Линейное: $T$ и её производные входят в первой степени. Однородное:
свободного члена нет, $T\equiv0$ — решение. Коэффициент $1/r$ переменный.

\q{Какая это задача: Коши, краевая или смешанная?}
Краевая: времени и начальных условий нет, условия только на границе.
Смешанной в лекциях называется начально-краевая задача для волнового
уравнения или уравнения теплопроводности. У нас «смешаны» только роды
граничных условий: I~род на боковой поверхности, II~род на торцах.

\q{Это задача Дирихле или Неймана?}
Ни то ни другое в чистом виде. Поэтому нет ни условия разрешимости задачи
Неймана, ни произвола в константе: уровень температуры зафиксирован
условием на боковой поверхности.

\q{Почему для уравнения Лапласа не ставят задачу Коши?}
Она некорректна: решение не зависит непрерывно от данных. Пример Адамара:
$u_{xx}+u_{yy}=0$, $u(x,0)=0$, $u_y(x,0)=\frac1n\sin nx$; решение
$u=\frac{1}{n^2}\sin nx\,\operatorname{sh}ny$ при $n\to\infty$ неограниченно
растёт, хотя данные стремятся к нулю.

\q{Что значит «задача поставлена корректно»?}
Решение существует, единственно и непрерывно зависит от данных (малое
изменение условий даёт малое изменение решения).

\q{Докажите единственность.}
Разность двух решений $w$ удовлетворяет задаче с нулевыми данными:
$\Delta w=0$, $w=0$ на боковой поверхности, $\partial w/\partial n=0$ на
торцах. По формуле Остроградского для вектора $w\,\Grad w$:
$\iiint_\Omega|\Grad w|^2dV=\iint_{S_\Omega} w\,\frac{\partial w}{\partial n}\,dS=0$.
Значит, $\Grad w=0$, $w$ постоянна, а на боковой поверхности $w=0$, поэтому
$w\equiv0$. Оговорка: из-за разрыва граничной функции строго это делается
в классе ограниченных решений (п.~\ref{der:uniq}).

\q{Что значит «единственность в классе ограниченных решений»? Бывают
неограниченные?}
Да. Пример — производная решения по параметру $c$: $v=\partial T/\partial c$.
Уравнение и условия на торцах можно дифференцировать по $c$, поэтому $v$ им
удовлетворяет; на боковой поверхности $v=0$ при всех $z\ne c$, потому что
там $\mu$ от $c$ не зависит; на оси $v$ ограничена. Но около точки $(a,c)$
она растёт как $1/d$, где $d$ — расстояние до точки. Функция $T+B\,v$
с любым числом $B$ удовлетворяет уравнению, условиям на торцах и на оси
и условию на боковой поверхности при всех $z\ne c$, но не ограничена
в области~$\Pi$. Поэтому в постановке отдельно требуется ограниченность
в $\Pi$: без неё решений бесконечно много.

\q{В каком классе функций ищется решение? Что такое $C^2(D)$ и
$C(\overline D)$?}
$C^2(D)$ — непрерывные вторые производные внутри области. $C(\overline D)$
— непрерывность вплоть до границы. У нас граничная функция разрывна,
поэтому вместо $C(\overline D)$ — ограниченность и непрерывность всюду,
кроме точки $(a,c)$.

\q{В лекции граница $\Gamma\in C^2$ — гладкая. У цилиндра граница гладкая?}
Нет, только кусочно-гладкая: на рёбрах (окружности $r=a$, $z=0$ и $r=a$,
$z=h$) нормаль не определена. Поэтому условия ставятся на гладких частях:
II~рода на торцах при $r<a$, I~рода на боковой поверхности при $0<z<h$; на
рёбрах ничего не задаётся. На решение это не влияет: рёбра имеют нулевую
площадь, тепло через них не проходит.

\q{Что такое разрыв первого рода?}
Точка, в которой у функции существуют конечные пределы слева и справа, но
они не равны. У $\mu(z)$ при $z=c$ предел слева равен $0$, справа — $T_0$,
скачок равен $T_0$. Функция с конечным числом таких разрывов
кусочно-непрерывна, и её можно разложить в ряд Фурье по косинусам
(лекция~5: ряд сходится в среднем квадратическом). В самой точке разрыва
сумма ряда равна полусумме пределов, $T_0/2$ (теорема Дирихле).

\q{В каких пределах лежит температура?}
$0\le T\le T_0$ по принципу максимума (при $T_0>0$). Экстремум не может
быть на торцах: после чётного продолжения по $z$ торцы становятся
внутренними точками.

\q{В отчёте сказано «пусть для определённости $T_0>0$». А если $T_0<0$?}
Постановка~\rf{eq:problem} не меняется: знак $T_0$ в ней нигде не
использован. Меняется физическая картина (верхняя часть холоднее нижней,
тепло идёт снизу вверх) и оценка~\rf{eq:bounds}: теперь $T_0\le T\le0$. В
общем виде $\min(0,T_0)\le T\le\max(0,T_0)$. Это видно и
из~\rf{eq:scaling}: $T=T_0\,\theta$, а $0\le\theta\le1$ при любом знаке
$T_0$.

\q{Почему наибольшее значение не может достигаться на оси? Ведь $r=0$ —
сторона прямоугольника $\Pi$.}
Принцип максимума применяется не к прямоугольнику, а к цилиндру: там $T$ —
гармоническая функция трёх переменных $x$, $y$, $z$. Точки оси лежат внутри
цилиндра, поэтому экстремума в них быть не может. Сторона $r=0$ появилась
только из-за перехода к координатам $(r,z)$.

\q{Существует ли решение?}
Да, оно строится явно методом Фурье на следующем этапе.

\q{Что будет с решением, если граничную температуру немного изменить?}
Изменится не больше, чем граничная температура:
$|T_1-T_2|\le\max|\mu_1-\mu_2|$ — следствие принципа максимума.

\q{Физичен ли скачок температуры на границе?}
Это идеализация узкой переходной зоны. Цена: в точке $(a,c)$ градиент
температуры и плотность потока неограниченны, хотя сама температура
ограничена. При сглаживании скачка (переходная зона ширины~$\delta$)
решение заметно меняется только на расстояниях порядка~$\delta$ от линии
скачка, а вдали от неё — мало. Оценка $|T_1-T_2|\le\max|\mu_1-\mu_2|$ из
предыдущего вопроса этого не показывает: при сглаживании
$\max|\mu_1-\mu_2|$ не мал (до $T_0/2$). Малость вдали от скачка следует из
того, что решение во внутренней точке выражается через интеграл от~$\mu$,
а $\int_0^h|\mu_1-\mu_2|\,dz\sim T_0\delta$ мал.

\q{Зачем безразмерная форма и сколько в задаче независимых параметров?}
Два: $a/h$ и $c/h$. $T_0$ — масштаб. Значит, графики достаточно строить
в единичном квадрате и исследовать влияние только двух чисел.

\q{Почему $r$ делится на $a$, а $z$ — на $h$? Можно ли взять один масштаб
длины?}
Можно. Если обе координаты делить на $a$, уравнение не будет содержать
параметров, зато область станет прямоугольником $0<\xi<1$, $0<\zeta<h/a$, а
скачок окажется при $\zeta=c/a$; параметров по-прежнему два. Выбранный
вариант удобен тем, что область — всегда единичный квадрат, а оба параметра
видны явно: $a/h$ в уравнении, $c/h$ в граничном условии.

\subsection{«А что, если…»}

\q{Если торцы не изолированы, а поддерживаются при заданной температуре?}
На торцах будут условия I~рода, получится задача Дирихле. При нулевой
температуре торцов при разделении переменных по $z$ вместо косинусов
появятся синусы $\sin(\pi nz/h)$. Если торцы имеют постоянные температуры
$T_{\text{н}}$ и $T_{\text{в}}$, сначала вычитают линейную функцию
$T_{\text{н}}+(T_{\text{в}}-T_{\text{н}})\,z/h$: она удовлетворяет уравнению
Лапласа, и для разности условия на торцах нулевые — снова синусы.

\q{Если на торцах теплообмен со средой?}
Условия III~рода. Закон Ньютона: тепло, уходящее с единицы площади
поверхности в среду, пропорционально разности температур поверхности и
среды, $-k\,\partial T/\partial n=\gamma\,(T-T_{\text{ср}})$. Здесь
$\gamma>0$ — коэффициент теплоотдачи, Вт/(м$^2\cdot$К), $T_{\text{ср}}$ —
температура среды. Обозначим $H=\gamma/k$ (размерность 1/м; к высоте $h$
эта буква отношения не имеет). На верхнем торце
$\partial T/\partial n=T_z$, на нижнем $\partial T/\partial n=-T_z$, поэтому
$T_z+H(T-T_{\text{ср}})=0$ при $z=h$ и $T_z-H(T-T_{\text{ср}})=0$ при $z=0$.
При $H\to0$ получается наша изоляция, при $H\to\infty$ — условие I~рода
$T=T_{\text{ср}}$.

\q{Если внутри есть источники тепла?}
Уравнение Пуассона $\Delta T=-F/k$ вместо уравнения Лапласа; граничные
условия те же.

\q{Если цилиндр полый, $b<r<a$?}
Ось не входит в область, условие ограниченности не нужно; вместо него
ставится граничное условие на внутренней поверхности $r=b$.

\q{Если изолирована вся поверхность?}
Задача Неймана с нулевым потоком: решение — любая постоянная, единственности
нет.

\q{Если нужен переходный процесс, а не установившийся режим?}
Уравнение $T_t=\chi\bigl(T_{rr}+\frac1r T_r+T_{zz}\bigr)$, те же граничные
условия плюс начальное $T|_{t=0}$. Это уже смешанная задача для уравнения
параболического типа. Наше решение — её предел при $t\to\infty$.

\q{Если теплопроводность зависит от координат?}
Вместо $\Delta T=0$ будет $\Div(k\,\Grad T)=0$ — $k$ нельзя вынести и
сократить.

\subsection{По списку источников и о следующем этапе}

\q{Какие источники использованы и что из каждого взято?}
В списке три источника. Первый — Куликов, Фирсов: учебное пособие СПбПУ по
уравнениям математической физики, первое в списке литературы на слайдах
лекций; общий источник по курсу. Второй — Тихонов, Самарский: закон Фурье и
вывод уравнения теплопроводности, классификация уравнений второго порядка
(гл.~I), уравнение Лапласа, принцип максимума, единственность, в том числе
при разрывных граничных условиях (гл.~IV). Третий — Смирнов, т.~2: вывод
уравнения теплопроводности с формулой Остроградского (п.~129), оператор
Лапласа в цилиндрических координатах (п.~131), свойства гармонических
функций и принцип максимума (п.~204).

\q{Каким методом будете решать и почему постановка это позволяет?}
Методом разделения переменных (Фурье): $T=R(r)Z(z)$. Уравнение линейное и
однородное, область — прямоугольник в координатах $(r,z)$, условия по $z$
однородные. По $z$ получится задача Штурма–Лиувилля $Z''+\lambda Z=0$,
$Z'(0)=Z'(h)=0$ с собственными функциями $\cos(\pi nz/h)$, $n=0,1,2,\dots$

\q{Почему собственные функции будут по $z$, а не по $r$?}
Задача Штурма–Лиувилля требует однородных граничных условий. Они заданы на
торцах, то есть по $z$. Неоднородное условие при $r=a$ используется
в конце — из него находятся коэффициенты ряда.

% ============================================================
\section{Ловушки}

\begin{enumerate}
  \item Назвать задачу «смешанной». В курсе смешанная — начально-краевая.
  \item Потерять $\frac1r T_r$ и написать $T_{rr}+T_{zz}=0$ — это плоская
  задача, а не цилиндр.
  \item Забыть условие на оси или назвать его граничным условием II~рода.
  \item Записать изоляцию как $T=0$ вместо $\partial T/\partial z=0$.
  \item Сказать «не зависит от $\varphi$, потому что очевидно». Нужны тело
  вращения, условия без $\varphi$ и единственность трёхмерной задачи.
  \item Перепутать $a$ (радиус) с $a^2$ из лекций, $c$ (высота скачка)
  с теплоёмкостью, $D$ (дискриминант) с областью, $a_{11}$ с радиусом.
  \item Сказать, что решение принадлежит $C(\overline D)$: граничная функция
  разрывна.
  \item Определять тип по уравнению, умноженному на $r$: получится
  $D=-r^2$, ноль на оси. Это особенность координат; в декартовых
  координатах уравнение Лапласа эллиптично всюду.
  \item Считать, что в ответ войдут $k$, $C$, $\rho$.
  \item Перепутать пояса: ноль — внизу ($z<c$), $T_0$ — вверху ($z\ge c$).
  \item Ждать произвольной константы или условия разрешимости, как в задаче
  Неймана.
  \item На нижнем торце сказать $\partial T/\partial n=T_z$ (там минус).
  \item Сказать, что $\ln r$ — источник: при $T=\ln r$ это сток (поток
  направлен к оси); в общем случае $C_1\ln r$ — источник или сток.
  \item Сказать «$r\,T_r$ сохраняется»: это верно только при $T_{zz}=0$.
\end{enumerate}

% ============================================================
\section{Самопроверка}

Могу ли я без подсказок:
\begin{itemize}[label=$\square$]
  \item за три минуты рассказать путь от условия до постановки~\rf{eq:problem};
  \item написать закон Фурье и баланс тепла, получить уравнение теплопроводности;
  \item объяснить переход к уравнению Лапласа;
  \item написать оператор Лапласа в цилиндрических координатах и вывести
  его осесимметричную часть;
  \item вывести основное уравнение из баланса тепла в кольце;
  \item объяснить каждое граничное условие: физика, формула, род, знак нормали;
  \item объяснить условие на оси и пример с $\ln r$;
  \item записать трёхмерную постановку~\rf{eq:problem3d} и объяснить
  условие периодичности;
  \item нарисовать рисунки~1 и~2 отчёта по памяти, подписать все стороны и
  объяснить штриховку, заливку, точку $(a,c)$ и нормали;
  \item определить тип уравнения через $D$, записать характеристическое
  уравнение;
  \item доказать единственность формулой Остроградского и назвать оговорку
  про разрыв;
  \item объяснить, почему из единственности следует осевая симметрия;
  \item объяснить $0\le T\le T_0$;
  \item перейти к безразмерным переменным и назвать два параметра;
  \item назвать каждую букву в любой формуле отчёта и её размерность.
\end{itemize}

\end{document}
